\documentclass[11pt,oneside]{article}
\usepackage[a4paper,margin=27mm,headheight=15pt]{geometry}
\usepackage[T1]{fontenc}
\usepackage[utf8]{inputenc}
\IfFileExists{libertinus.sty}{\usepackage{libertinus}}{\usepackage{lmodern}}
\usepackage{microtype}
\usepackage{amsmath,amssymb,amsthm,mathtools,mathrsfs,bm}
\usepackage{booktabs,longtable,array,tabularx}
\usepackage{xcolor}
\usepackage{enumitem}
\usepackage{fancyhdr}
\usepackage{titlesec}
\usepackage{listings}
\usepackage{xurl}
\usepackage{hyperref}
\usepackage[nameinlink,noabbrev]{cleveref}

\definecolor{linkblue}{RGB}{25,75,135}
\definecolor{shadegray}{RGB}{246,247,249}
\definecolor{rulegray}{RGB}{195,201,208}
\hypersetup{
  colorlinks=true,
  linkcolor=linkblue,
  citecolor=linkblue,
  urlcolor=linkblue,
  pdftitle={The exact Fourier-decay theory of Rvachev's up-function},
  pdfsubject={Rigorous reconciliation of four proposed solutions, the lacunary-product literature, and the Cesaro normalization problem},
  pdfkeywords={Rvachev up-function, Fabius function, Fourier transform, sinc product, transfer operator, Thue-Morse, Cesaro mean, law of the iterated logarithm, Lambert W}
}

\pagestyle{fancy}
\fancyhf{}
\fancyhead[L]{\small Rvachev up-function}
\fancyhead[R]{\small Final Fourier-decay synthesis}
\fancyfoot[C]{\thepage}
\renewcommand{\headrulewidth}{0.4pt}
\titleformat{\section}{\Large\bfseries}{\thesection}{0.7em}{}
\titleformat{\subsection}{\large\bfseries}{\thesubsection}{0.7em}{}
\titleformat{\subsubsection}{\normalsize\bfseries}{\thesubsubsection}{0.7em}{}
\setcounter{secnumdepth}{3}
\setcounter{tocdepth}{2}
\setlist[itemize]{leftmargin=2em,itemsep=0.25em,topsep=0.4em}
\setlist[enumerate]{leftmargin=2.3em,itemsep=0.25em,topsep=0.4em}
\setlength{\emergencystretch}{3em}
\allowdisplaybreaks

\newtheorem{theorem}{Theorem}[section]
\newtheorem{proposition}[theorem]{Proposition}
\newtheorem{lemma}[theorem]{Lemma}
\newtheorem{corollary}[theorem]{Corollary}
\newtheorem{conjecture}[theorem]{Conjecture}
\theoremstyle{definition}
\newtheorem{definition}[theorem]{Definition}
\newtheorem{construction}[theorem]{Construction}
\newtheorem{example}[theorem]{Example}
\theoremstyle{remark}
\newtheorem{remark}[theorem]{Remark}
\newtheorem{warning}[theorem]{Warning}
\newtheorem{auditfinding}[theorem]{Audit finding}
\crefname{auditfinding}{audit finding}{audit findings}
\Crefname{auditfinding}{Audit finding}{Audit findings}

\newcommand{\R}{\mathbb R}
\newcommand{\C}{\mathbb C}
\newcommand{\N}{\mathbb N}
\newcommand{\Z}{\mathbb Z}
\newcommand{\e}{\mathrm e}
\newcommand{\ii}{\mathrm i}
\newcommand{\dd}{\,\mathrm d}
\newcommand{\Up}{\operatorname{up}}
\newcommand{\sinc}{\operatorname{sinc}}
\newcommand{\bigO}{\mathcal O}
\newcommand{\smallo}{o}
\newcommand{\supp}{\operatorname{supp}}
\newcommand{\sgn}{\operatorname{sgn}}
\newcommand{\Var}{\operatorname{Var}}
\newcommand{\E}{\mathbb E}
\newcommand{\Prob}{\mathbb P}
\newcommand{\Lip}{\operatorname{Lip}}
\newcommand{\vtwo}{\nu_2}
\newcommand{\Lop}{\mathcal L}
\newcommand{\T}{T}
\newcommand{\Wlam}{W_0}
\newcommand{\abs}[1]{\left|#1\right|}
\newcommand{\norm}[1]{\left\lVert#1\right\rVert}
\newcommand{\set}[1]{\left\{#1\right\}}
\newcommand{\floor}[1]{\left\lfloor#1\right\rfloor}
\newcommand{\ceil}[1]{\left\lceil#1\right\rceil}

\newenvironment{boxedremark}{%
  \par\smallskip\noindent\begin{center}\begin{minipage}{0.94\textwidth}%
  \color{black}\hrule\vspace{0.6em}\small
}{%
  \vspace{0.6em}\hrule\end{minipage}\end{center}\smallskip
}

\lstdefinestyle{code}{
  basicstyle=\ttfamily\footnotesize,
  backgroundcolor=\color{shadegray},
  frame=single,
  rulecolor=\color{rulegray},
  columns=fullflexible,
  keepspaces=true,
  showstringspaces=false,
  breaklines=true,
  xleftmargin=0.7em,
  xrightmargin=0.7em,
  aboveskip=0.8em,
  belowskip=0.8em
}

\title{\bfseries The Exact Fourier-Decay Theory of Rvachev's Up-Function\\[0.35em]
\large Reconciliation of four proposed solutions, correction of the lacunary-product LIL,\\
and a complete resolution of the Ces\`aro normalization problem}
\author{Final synthesis for Vladimir Reshetnikov's \texttt{ProveIt} project}
\date{26 August 2026}

\begin{document}
\maketitle

\begin{abstract}
Let
\[
 \Phi(x)=\widehat{\Up}(x)
 =\prod_{n=0}^{\infty}\sinc\!\left(\frac{\pi x}{2^n}\right),
 \qquad \sinc u=\frac{\sin u}{u},
\]
be the Fourier transform of Rvachev's compactly supported up-function, with
Fourier convention \(\widehat h(\xi)=\int_{\R}h(t)\e^{-2\pi\ii t\xi}\dd t\).
The product is rapidly decreasing but has zeros at every nonzero integer and
contains a lacunary dynamical factor.  Consequently there is no single
pointwise ``decay exponent'': typical values, arithmetic means, quadratic
means, fixed-ray extremals, dyadic-annulus maxima, and two-adically exceptional
lobes have genuinely different lower-order behavior.

This article audits the four wave--1 analyses, the paper of
Aistleitner--Hofer--Larcher, and the wave--2 comparative report, and then closes
the remaining logical gap.  The exact dyadic factorization reduces the problem
to
\(P_n(z)=\prod_{j=0}^{n-1}|\sin(\pi2^jz)|\).  Its arithmetic mean is governed by
the Perron root
\(\rho=0.661322602060565\ldots\) (numerical; rigorously
\(0.66130<\rho<0.66135\)) of an explicit transfer operator.  Thus the natural
arithmetic gauge is
\[
 g_1(x)=\exp\!\left(-\frac{(\log x)^2}{2\log2}\right)x^{-\kappa_1},
 \qquad
 \kappa_1=\frac12+\log_2\frac{\pi}{\rho}
 =2.748070014871334\ldots .
\]
After multiplication by
\(A_1=0.09126612413159\ldots\), this gauge answers the bounded form of the
original question and gives limit one along dyadic endpoints.  At unrestricted
times its means approach a nonconstant continuous mantissa profile generated by
a singular transfer-operator eigenmeasure.  This proves impossibility for every
gauge whose shell distortion relative to \(g_1\) is bounded, or whose shell
shape stabilizes; it includes all Hardy-field gauges of the correct scale.

The principal new theorem is that this obstruction is not an unrestricted
nonexistence theorem.  By equalizing the mass on successively finer dyadic
partitions inside each shell, respecting the one-sided derivative constraint
imposed by monotonicity, and then applying the Hoischen--Carleman theorem on
simultaneous entire approximation of a function and its derivative, we
construct a positive real-analytic, eventually strictly decreasing gauge
\(g\) for which
\[
 \lim_{t\to\infty}\frac1{t-t_0}\int_{t_0}^t\frac{|\Phi(x)|}{g(x)}\dd x=1
\]
for every fixed \(t_0\).  It can be chosen with
\(\log g(x)=-(\log x)^2/(2\log2)-\kappa_1\log x+o(\log x)\).
Any such exact gauge must, however, develop unbounded and increasingly fine
within-shell distortion; no simple stabilized formula can do the job.

We also give the complete exponent spectrum
\(\kappa_0=3.15149612947\ldots\) (typical),
\(\kappa_1=2.74807001487\ldots\) (arithmetic),
\(\kappa_2=\log_2(2\pi)=2.65149612947\ldots\) (RMS), and
\(\kappa_\infty=2.35901487911\ldots\) (uniform coarse envelope), prove the
sharp Lambert-\(W\) asymptotic for the largest peak in a dyadic annulus, and
rederive the central limit theorem and law of the iterated logarithm.  The exact
variance is \(\pi^2/4\), not \(\pi^2/2\); hence the published LIL constant in
the cited lacunary-product paper is too large by \(\sqrt2\).
\end{abstract}

\begin{boxedremark}
\textbf{Final answer in one paragraph.}
There is no universal pointwise equivalent for \(|\Phi(x)|\).  The simple
analytic gauge \(A_1g_1\) gives the correct arithmetic scale, satisfies the
two-sided bounded requirement, and has a nonconstant dyadic-mantissa Ces\`aro
profile.  Every bounded-distortion, stabilizing, or Hardy-field modification of
that gauge fails to make the unrestricted Ces\`aro limit equal to one.  However,
the literal existence question has an affirmative answer: an adaptive
non-elementary positive analytic decreasing gauge can be constructed so that
the unrestricted Ces\`aro limit is exactly one.  Its increasingly fine shell
microstructure is not optional; it is precisely what escapes the singular
measure obstruction.
\end{boxedremark}

\tableofcontents
\newpage

\section{Problem, conventions, and the hierarchy of answers}
\label{sec:introduction}

The original question asks for a real-valued gauge \(g\), eventually analytic,
positive, and monotone decreasing, such that either
\begin{equation}
 \frac1{t-t_0}\int_{t_0}^t\frac{|\Phi(x)|}{g(x)}\dd x\longrightarrow1
 \qquad(t\to\infty),
 \label{eq:Q2}
\end{equation}
or at least
\begin{equation}
 0<A<\frac1{t-t_0}\int_{t_0}^t\frac{|\Phi(x)|}{g(x)}\dd x<B<\infty
 \qquad(t\ge t_1).
 \label{eq:Q3}
\end{equation}
The phrase ``decay rate'' is ambiguous before one specifies what is being
averaged.  Four elementary facts already force this distinction.

\begin{enumerate}
\item \(\Phi(m)=0\) for every nonzero integer \(m\), with nonuniform
multiplicity.  Therefore no positive pointwise equivalent can hold on the full
real axis.
\item On a fixed dyadic ray \(x=2^ky\), the remaining oscillation is a Birkhoff
sum for the doubling map.  Typical rays fluctuate on a Gaussian and
iterated-logarithm scale.
\item The arithmetic and quadratic shell normalizations are governed by
different transfer-operator eigenvalues.
\item The largest peak in \([2^k,2^{k+1}]\) is a boundary-layer optimization:
it begins near the zero at \(2^k\), then shadows the maximizing orbit
\(1/3\leftrightarrow2/3\).  This produces a Lambert-\(W\) correction absent
from fixed-ray asymptotics.
\end{enumerate}

Write throughout
\begin{equation}
 \lambda=\log2,
 \qquad
 E(x)=\exp\!\left(-\frac{(\log x)^2}{2\lambda}\right).
 \label{eq:lambda-E}
\end{equation}
The following table gives the correct power correction attached to the common
log-normal factor \(E(x)\).  The precise meaning of each row is developed below.

\begin{center}
\begin{tabularx}{0.96\textwidth}{@{}>{\raggedright\arraybackslash}Xll@{}}
\toprule
notion of normalized amplitude & product scale & power exponent \(\kappa\)\\
\midrule
Lebesgue-typical dyadic ray / geometric mean
 & \(\Lambda_0=\tfrac12\)
 & \(\kappa_0=\tfrac32+\log_2\pi=3.151496129472319\ldots\)\\[0.25em]
arithmetic normalization
 & \(\Lambda_1=\rho\)
 & \(\kappa_1=\tfrac12+\log_2(\pi/\rho)=2.748070014871334\ldots\)\\[0.25em]
root-mean-square normalization
 & \(\Lambda_2=2^{-1/2}\)
 & \(\kappa_2=\log_2(2\pi)=2.651496129472318\ldots\)\\[0.25em]
uniform coarse envelope
 & \(\Lambda_\infty=\sqrt3/2\)
 & \(\kappa_\infty=\tfrac32+\log_2(\pi/\sqrt3)=2.359014879111741\ldots\)\\
\bottomrule
\end{tabularx}
\end{center}

The ordering
\[
 \kappa_\infty<\kappa_2<\kappa_1<\kappa_0
\]
is the usual ordering of extremal, quadratic, arithmetic, and geometric
behavior.  It is not a disagreement among valid calculations; it is a spectrum
of answers to inequivalent questions.

\section{Exact dyadic factorization and zero geometry}
\label{sec:factorization}

\subsection{The shell identity}

For \(x>0\), write uniquely
\[
 x=2^ky,
 \qquad k\in\Z,\quad 1\le y<2,
 \qquad z=y-1\in[0,1).
\]
Define
\begin{equation}
 P_n(z)=\prod_{j=0}^{n-1}\abs{\sin(\pi2^jz)},
 \qquad
 H(y)=\prod_{m=1}^{\infty}\abs{\sinc\!\left(\frac{\pi y}{2^m}\right)}.
 \label{eq:P-H}
\end{equation}
The product \(H\) is continuous on \([1,2]\), positive on \([1,2)\), and
vanishes at \(2\).

\begin{proposition}[exact dyadic-shell factorization]
\label{prop:factorization}
For every integer \(k\ge0\), \(1\le y\le2\), and \(n=k+1\),
\begin{equation}
 \boxed{
 |\Phi(2^ky)|
 =2^{-k(k+1)/2}(\pi y)^{-n}H(y)P_n(y-1).}
 \label{eq:factorization}
\end{equation}
\end{proposition}

\begin{proof}
Separate the factors with indices \(0\le m\le k\) from the tail.  In the
finite part put \(j=k-m\).  Since \(2^j\) is an integer,
\[
 \abs{\sin(\pi2^jy)}=\abs{\sin(\pi2^j(y-1))}.
\]
The finite denominators contribute
\[
 \prod_{j=0}^k \pi2^jy=(\pi y)^{k+1}2^{k(k+1)/2},
\]
and the remaining factors are exactly \(H(y)\).  Multiplication gives
\eqref{eq:factorization}.
\end{proof}

The universal quadratic term is already visible in
\(2^{-k(k+1)/2}\).  Everything delicate is contained in \(P_n\), together
with the competition between its orbit and the factor \(y^{-n}\).

\subsection{Canonical product, multiplicities, and one peak per lobe}

The Euler product for \(\sinc\) can be regrouped without convergence
ambiguity because the resulting genus-zero product is absolutely convergent.

\begin{proposition}[canonical zero product]
\label{prop:canonical}
For every \(z\in\C\),
\begin{equation}
 \Phi(z)=\prod_{m=1}^{\infty}
 \left(1-\frac{z^2}{m^2}\right)^{1+\vtwo(m)}.
 \label{eq:canonical}
\end{equation}
Thus the zeros at \(\pm m\) have multiplicity \(1+\vtwo(m)\).
\end{proposition}

\begin{proof}
Expand each factor as
\[
 \sinc\!\left(\frac{\pi z}{2^h}\right)
 =\prod_{r=1}^{\infty}\left(1-\frac{z^2}{r^2 4^h}\right).
\]
A positive integer \(m\) is representable as \(r2^h\) exactly for
\(0\le h\le\vtwo(m)\), giving the stated multiplicity.  Finally,
\(\sum_m(1+\vtwo(m))/m^2<\infty\), so the regrouping is legitimate.
\end{proof}

\begin{theorem}[strict log-concavity between consecutive zeros]
\label{thm:one-peak}
For every integer \(m\ge1\), \(|\Phi|\) has exactly one critical point in
\((m,m+1)\); it is the unique maximum of that lobe.
\end{theorem}

\begin{proof}
Away from the zeros, termwise differentiation in \eqref{eq:canonical} gives
\begin{equation}
 \frac{\dd^2}{\dd x^2}\log|\Phi(x)|
 =-2\sum_{r=1}^{\infty}(1+\vtwo(r))
 \frac{r^2+x^2}{(r^2-x^2)^2}<0.
 \label{eq:strict-log-concavity}
\end{equation}
Hence the logarithmic derivative is strictly decreasing on every interzero
interval.  It tends to \(+\infty\) at the left endpoint and \(-\infty\) at
the right endpoint, so it has one and only one zero.
\end{proof}

This theorem is useful both conceptually and computationally: the lobe maxima
can be located by one-dimensional root finding, and exceptional multiplicities
can be analyzed locally.

\section{Transfer operator and the arithmetic eigenvalue}
\label{sec:transfer}

\subsection{The doubling map and its weighted Perron operator}

Let
\[
 \T z=2z\pmod1,
 \qquad w(z)=|\sin\pi z|.
\]
Then \(P_n(z)=\prod_{j=0}^{n-1}w(\T^jz)\).  For \(q>0\), define
\begin{equation}
 (\Lop_qh)(x)=\frac12\left[
  \sin^q\!\left(\frac{\pi x}{2}\right)h\!\left(\frac x2\right)
 +\cos^q\!\left(\frac{\pi x}{2}\right)h\!\left(\frac{x+1}{2}\right)
 \right].
 \label{eq:Lq}
\end{equation}
A change of variables on the two inverse branches gives
\begin{equation}
 \int_0^1 h(z)P_n(z)^q\dd z
 =\int_0^1\Lop_q^n h(x)\dd x.
 \label{eq:transfer-identity}
\end{equation}
We write \(\Lop=\Lop_1\).

\begin{theorem}[Ruelle--Perron--Frobenius package for \(q=1\)]
\label{thm:RPF}
On the Banach space of Lipschitz functions on \([0,1]\), the operator
\(\Lop\) has a simple positive leading eigenvalue \(\rho\), a strictly
positive Lipschitz eigenfunction \(h\), and a positive eigenprobability
\(\nu\), normalized by \(\nu(h)=1\), such that
\begin{equation}
 \rho^{-n}\Lop^n\varphi
 \longrightarrow h\,\nu(\varphi)
 \quad\text{uniformly on }[0,1]
 \label{eq:RPF-limit}
\end{equation}
for every Lipschitz \(\varphi\).  Moreover
\begin{equation}
 \rho=0.661322602060565\ldots\quad\text{numerically},
 \qquad 0.66130<\rho<0.66135\quad\text{rigorously}.
 \label{eq:rho}
\end{equation}
\end{theorem}

\begin{proof}[Proof architecture]
The elementary Lasota--Yorke estimate has the form
\[
 \Lip(\Lop\varphi)
 \le \frac{\sqrt2}{4}\Lip(\varphi)+C\norm{\varphi}_\infty,
\]
while \(\norm{\Lop\varphi}_\infty\le2^{-1/2}\norm{\varphi}_\infty\).
Thus the essential spectral radius on the Lipschitz space is at most
\(\sqrt2/4\).  The exact second moment \(\int_0^1P_n^2=2^{-n}\), together with the
Gelfond estimate proved in \cref{sec:extremal}, gives
\[
 \int_0^1P_n\,\dd z
 \ge \frac{\int_0^1P_n^2\,\dd z}{\lVert P_n\rVert_\infty}
 \ge \sqrt{\frac35}\left(\frac1{\sqrt3}\right)^n.
\]
Hence \(\rho\ge1/\sqrt3>\sqrt2/4\), so the leading spectrum is isolated.
Positivity, topological exactness of the doubling map, and the fact that the
weights vanish only at branch endpoints give irreducibility and simplicity.
This yields \eqref{eq:RPF-limit} by the standard positive
quasi-compact-operator argument.  The rigorous enclosure in
\eqref{eq:rho} is the one proved in the lacunary-product paper by iterating
concave upper and lower ratios; the longer decimal is a stable numerical
collocation value.
\end{proof}

For later use set
\begin{equation}
 H_0=\int_0^1h(x)\dd x.
 \label{eq:H0}
\end{equation}
Combining \eqref{eq:transfer-identity} and \eqref{eq:RPF-limit} gives
\begin{equation}
 \rho^{-n}\int_0^1\varphi(z)P_n(z)\dd z
 \longrightarrow H_0\nu(\varphi)
 \label{eq:weighted-P-limit}
\end{equation}
for continuous \(\varphi\); the extension from Lipschitz to continuous
functions follows by uniform approximation.

\subsection{The scale dictionary}

Suppose a product norm grows like \(\Lambda^n\).  The power correction needed
to cancel it in the exact shell factorization is
\begin{equation}
 \boxed{
 \kappa(\Lambda)=
 \frac{\log\pi+\lambda/2-\log\Lambda}{\lambda}.}
 \label{eq:kappa-dictionary}
\end{equation}
Indeed, for
\[
 g_\kappa(x)=E(x)x^{-\kappa},
\]
a direct calculation from \eqref{eq:factorization} gives
\begin{equation}
 \frac{|\Phi(2^k(1+z))|}{g_\kappa(2^k(1+z))}
 =\Lambda^{-n}W_\Lambda(z)P_n(z),
 \qquad n=k+1,
 \label{eq:normalized-shell-general}
\end{equation}
where
\begin{equation}
 W_\Lambda(z)=\frac{\Lambda}{\pi}H(1+z)(1+z)^{\kappa-1}
 \exp\!\left(\frac{\log^2(1+z)}{2\lambda}\right).
 \label{eq:Wlambda}
\end{equation}
This identity, not a heuristic averaging argument, is the common source of all
four exponents.

For the arithmetic problem, put
\begin{equation}
 \kappa_1=\kappa(\rho)
 =\frac12+\log_2\frac\pi\rho,
 \qquad
 g_1(x)=E(x)x^{-\kappa_1},
 \label{eq:g1}
\end{equation}
and abbreviate
\begin{equation}
 W(z)=W_\rho(z)
 =\frac\rho\pi H(1+z)(1+z)^{\kappa_1-1}
 \exp\!\left(\frac{\log^2(1+z)}{2\lambda}\right).
 \label{eq:W}
\end{equation}
Then
\begin{equation}
 \boxed{
 \frac{|\Phi(2^k(1+z))|}{g_1(2^k(1+z))}
 =\rho^{-n}W(z)P_n(z).}
 \label{eq:normalized-shell}
\end{equation}

\section{The natural arithmetic gauge and its singular phase profile}
\label{sec:natural-gauge}

\subsection{The shell constant and limiting measure}

Define the finite measures
\begin{equation}
 \alpha_k(\dd z)
 =\rho^{-(k+1)}W(z)P_{k+1}(z)\dd z,
 \qquad 0\le z\le1.
 \label{eq:alpha-k}
\end{equation}
By \eqref{eq:weighted-P-limit},
\begin{equation}
 \alpha_k\Longrightarrow\alpha,
 \qquad
 \alpha(\dd z)=H_0W(z)\nu(\dd z).
 \label{eq:alpha-limit}
\end{equation}
Set
\begin{equation}
 A_1=\alpha([0,1])=H_0\nu(W)
 =0.09126612413159\ldots\quad\text{(numerical)},
 \label{eq:A1}
\end{equation}
and normalize
\begin{equation}
 \mu(\dd z)=\frac{W(z)\nu(\dd z)}{\nu(W)},
 \qquad
 F(z)=\mu([0,z]).
 \label{eq:mu-F}
\end{equation}

The following regularity fact is crucial.

\begin{theorem}[the arithmetic eigenmeasure is continuous, singular, and full support]
\label{thm:nu-singular}
The eigenmeasure \(\nu\), and hence \(\mu\), has no atoms, is singular with
respect to Lebesgue measure, and gives positive mass to every nonempty open
subinterval of \((0,1)\).
\end{theorem}

\begin{proof}
For an atom at \(z\), the dual eigenrelation gives
\begin{equation}
 \rho\nu\{z\}=\frac12\sin(\pi z)\nu\{\T z\}.
 \label{eq:atom-relation}
\end{equation}
By \cref{prop:gelfond} below,
\(\rho\ge1/\sqrt3>1/2\).  Iterating
\eqref{eq:atom-relation} therefore yields
\[
 \nu\{z\}\le(2\rho)^{-n}\longrightarrow0,
\]
so \(\nu\) is atomless.

Suppose an absolutely continuous eigencomponent exists, with density
\(r\in L^1\), \(r\ge0\).  A change of variables in the dual relation gives
\begin{equation}
 \sin(\pi z)r(\T z)=\rho r(z),
 \qquad
 r(\T^nz)=\frac{\rho^n}{P_n(z)}r(z).
 \label{eq:density-eigen}
\end{equation}
Birkhoff's theorem and
\(\int_0^1\log|\sin\pi z|\dd z=-\log2\) imply
\[
 P_n(z)=\exp[-n(\log2+o(1))]
\]
for almost every \(z\).  On the set \(\{r>0\}\), the right side of
\eqref{eq:density-eigen} therefore grows as
\(\exp[n(\log(2\rho)+o(1))]\).  On the other hand, for every
\(0<\varepsilon<\log(2\rho)\), Markov's inequality and invariance of
Lebesgue measure give
\[
 \sum_{n\ge1}\operatorname{Leb}\{z:r(\T^nz)>\e^{\varepsilon n}\}
 =\sum_{n\ge1}\operatorname{Leb}\{r>\e^{\varepsilon n}\}<\infty.
\]
Borel--Cantelli contradicts the preceding exponential lower growth.  The
dual operator sends absolutely continuous measures to absolutely continuous
measures and singular measures to singular measures, so uniqueness of the
Lebesgue decomposition shows that the absolutely continuous part of \(\nu\)
is zero.

Finally, every nonempty open interval contains a dyadic cylinder on which some
iterate of \(\T\) maps onto \((0,1)\), with positive product weight away from
finitely many endpoints.  Applying the dual eigenrelation and using
atomlessness shows that its \(\nu\)-mass is positive.  Since \(W>0\) on
\([0,1)\), the same is true for \(\mu\).
\end{proof}

\subsection{Exact dyadic-mantissa Ces\`aro profile}

Let
\begin{equation}
 g_{\mathrm{nat}}(x)=A_1g_1(x).
 \label{eq:gnat}
\end{equation}
The natural gauge is positive and real-analytic on \((0,\infty)\), and is
strictly decreasing for all sufficiently large \(x\).

\begin{theorem}[natural-gauge profile]
\label{thm:profile}
Fix \(t_0>0\).  For every \(1\le y\le2\),
\begin{equation}
 \lim_{k\to\infty}
 \frac1{2^ky-t_0}\int_{t_0}^{2^ky}
 \frac{|\Phi(x)|}{g_{\mathrm{nat}}(x)}\dd x
 =\mathscr P(y),
 \label{eq:profile-limit}
\end{equation}
where
\begin{equation}
 \boxed{
 \mathscr P(y)=\frac{1+F(y-1)}{y}.}
 \label{eq:profile}
\end{equation}
The convergence is uniform in \(y\in[1,2]\).  The function \(\mathscr P\)
is continuous, satisfies \(\mathscr P(1)=\mathscr P(2)=1\), and is not
identically one.
\end{theorem}

\begin{proof}
By \eqref{eq:normalized-shell}, for \(0\le z\le1\),
\begin{align}
 \int_{2^k}^{2^k(1+z)}\frac{|\Phi(x)|}{A_1g_1(x)}\dd x
 &=2^k\frac{\alpha_k([0,z])}{A_1}.
 \label{eq:partial-shell}
\end{align}
Weak convergence to the atomless measure \(\alpha\) implies uniform
convergence of the corresponding distribution functions, hence
\[
 \frac{\alpha_k([0,z])}{A_1}\longrightarrow F(z)
\]
uniformly.  A full shell contributes \(2^k(1+o(1))\).  Summing all preceding
shells and using the geometric weights gives
\[
 \int_{t_0}^{2^k(1+z)}\frac{|\Phi|}{A_1g_1}
 =2^k\bigl(1+F(z)+o(1)\bigr),
\]
uniformly in \(z\).  Division by \(2^k(1+z)-t_0\) proves
\eqref{eq:profile}.  If \(\mathscr P\equiv1\), then
\(F(z)=z\), contradicting singularity of \(\mu\).
\end{proof}

\begin{corollary}[answer to the bounded question]
\label{cor:Q3}
The gauge \(g_{\mathrm{nat}}\) satisfies \eqref{eq:Q3}.  It also satisfies
\eqref{eq:Q2} along the dyadic sequence \(t=2^k\), but not for unrestricted
\(t\).
\end{corollary}

\begin{proof}
The continuous function \(\mathscr P\) has a positive minimum and finite
maximum on \([1,2]\), and the convergence in \cref{thm:profile} is uniform.
At \(y=1\) the limit is one.  Nonconstancy gives failure at some other
mantissa.
\end{proof}

A crude theorem-level bound follows without numerics: since \(0\le F\le1\),
\(1/2\le\mathscr P(y)\le2\).  Thus, after enlarging \(t_0\), one may for
example take \(A=1/3\), \(B=3\) in \eqref{eq:Q3}.

\subsection{Eventual behavior of the derivatives}

For the regularity requested in the original question, the natural gauge has
a useful fixed-order property.

\begin{proposition}[fixed-order alternating derivatives]
\label{prop:derivatives}
For every fixed integer \(m\ge0\),
\begin{equation}
 g_1^{(m)}(x)=x^{-m}g_1(x)Q_m(\log x),
 \label{eq:derivative-polynomial}
\end{equation}
where \(Q_m\) is a degree-\(m\) polynomial with leading term
\((-1)^m\lambda^{-m}X^m\).  Consequently
\((-1)^mg_1^{(m)}(x)>0\) for all sufficiently large \(x\), and each fixed
derivative is eventually monotone.
\end{proposition}

\begin{proof}
Apply \(\dd/\dd x=x^{-1}\dd/\dd(\log x)\) inductively to
\(g_1(x)=\exp[-L^2/(2\lambda)-\kappa_1L]\), \(L=\log x\).  The leading
term is multiplied at each step by \(-L/\lambda\).  Monotonicity of the
\(m\)-th derivative follows from the eventual sign of the \((m+1)\)-st.
\end{proof}

The threshold depends on \(m\); this is not a claim of complete monotonicity
on one common half-line.

\section{What the singular profile does and does not forbid}
\label{sec:obstructions}

The wave--2 audit correctly identified a strong obstruction, but stated its
scope too broadly.  We first formulate the exact obstruction, strengthen it,
and then explain why it still leaves room for the adaptive construction in
\cref{sec:adaptive}.

\subsection{Linearity forced by an unrestricted Ces\`aro limit}

\begin{lemma}[shell linearity]
\label{lem:shell-linearity}
If a positive gauge \(g\) satisfies \eqref{eq:Q2}, then for every
\(z\in[0,1]\), uniformly in \(z\),
\begin{equation}
 2^{-k}\int_{2^k}^{2^k(1+z)}\frac{|\Phi(x)|}{g(x)}\dd x
 \longrightarrow z.
 \label{eq:shell-linearity}
\end{equation}
In particular, the normalized measures of \(|\Phi|/g\) inside every dyadic
shell converge weakly to Lebesgue measure on \([0,1]\).
\end{lemma}

\begin{proof}
Apply \eqref{eq:Q2} at \(2^k\) and at \(2^k(1+z)\), subtract the two
asymptotics, and divide by \(2^k\).  Uniformity follows from monotonicity of
the partial integrals together with convergence on a fine fixed mesh.
\end{proof}

\subsection{Stabilizing and bounded-distortion gauges}

\begin{theorem}[stabilizing-shape obstruction]
\label{thm:stabilizing-obstruction}
Suppose there exist constants \(c_k>0\) and a continuous function
\(u:[0,1]\to(0,\infty)\) such that
\begin{equation}
 \frac{g(2^k(1+z))}{c_kg_1(2^k(1+z))}\longrightarrow u(z)
 \quad\text{uniformly in }z.
 \label{eq:stabilizing}
\end{equation}
Then \(g\) cannot satisfy \eqref{eq:Q2}.
\end{theorem}

\begin{proof}
Up to a scalar normalization, the shell measure for \(|\Phi|/g\) is
\[
 u(z)^{-1}\alpha_k(\dd z)(1+o(1)).
\]
Every weak limit is therefore proportional to
\(u^{-1}\alpha\), which is singular.  This contradicts the Lebesgue limit
forced by \cref{lem:shell-linearity}.
\end{proof}

The next statement removes the existence of any limiting shape.

\begin{theorem}[bounded-distortion obstruction]
\label{thm:bounded-distortion}
Suppose there are constants \(0<m<M<\infty\) and numbers \(c_k>0\) such
that, for all sufficiently large \(k\) and all \(z\in[0,1]\),
\begin{equation}
 m\le
 \frac{g(2^k(1+z))}{c_kg_1(2^k(1+z))}
 \le M.
 \label{eq:bounded-distortion}
\end{equation}
Then \(g\) cannot satisfy \eqref{eq:Q2}.
\end{theorem}

\begin{proof}
After cancelling the irrelevant scalar \(c_k\), write the normalized shell
measure as
\[
 \eta_k(\dd z)=\frac{q_k(z)\alpha_k(\dd z)}{\int q_k\dd\alpha_k},
 \qquad M^{-1}\le q_k\le m^{-1}.
\]
Thus \(\eta_k\le C\alpha_k\) in the sense that
\(\int\varphi\dd\eta_k\le C\int\varphi\dd\alpha_k\) for every
continuous \(\varphi\ge0\), with \(C\) independent of \(k\).  Any weak
limit \(\eta\) is therefore dominated by \(C\alpha\), hence singular.
Again this contradicts \cref{lem:shell-linearity}.
\end{proof}

\begin{corollary}[necessary shell microstructure]
\label{cor:necessary-microstructure}
Every gauge satisfying the unrestricted limit \eqref{eq:Q2} must have
unbounded shell distortion relative to \(g_1\), even after multiplication by
an arbitrary scalar \(c_k\) on each shell.
\end{corollary}

\subsection{Hardy fields and the overreach about ``elementary'' functions}

If \(g/g_1\) belongs to a Hardy field, put
\(h=\log(g/g_1)\).  Then \(xh'(x)\) has a limit in
\([ -\infty,+\infty]\).  A finite limit \(\gamma\) gives
\[
 h(2^k(1+z))-h(2^k)\longrightarrow\gamma\log(1+z)
\]
uniformly, so \cref{thm:stabilizing-obstruction} applies.  Infinite limits
force shell mass to concentrate at one endpoint, contradicting
\eqref{eq:shell-linearity}.  Hence no Hardy-field gauge of the correct rough
scale can solve \eqref{eq:Q2}.

It does \emph{not} follow that every elementary closed-form expression is
covered.  Trigonometric composition already leaves the Hardy-field world.
For example, for \(2c\varepsilon<1/\lambda\),
\begin{equation}
 g_{\varepsilon,c}(x)
 =g_1(x)\exp\!\bigl(\varepsilon\sin(c(\log x)^2)\bigr)
 \label{eq:oscillatory-elementary}
\end{equation}
is elementary, positive, analytic, and eventually decreasing, but its shell
shape does not stabilize.  It still fails \eqref{eq:Q2}, now by
\cref{thm:bounded-distortion}, because its distortion is uniformly bounded.
Thus the correct conclusion is:

\begin{boxedremark}
Hardy-field, stabilizing, and bounded-distortion gauges are impossible.  A
blanket statement about all elementary formulas is not justified.  The exact
gauge constructed below escapes by developing unbounded, finer-and-finer shell
structure.
\end{boxedremark}

\section{An analytic decreasing gauge with the exact unrestricted Ces\`aro limit}
\label{sec:adaptive}

This section resolves the literal existence question left open by all supplied
materials.  The proof has two stages.  First we construct a smooth decreasing
gauge by shellwise mass equalization.  Then we approximate its logarithm and
first derivative by an entire function with an arbitrarily small variable
error.

\subsection{A shellwise equalization lemma}

Recall the measures \(\alpha_k\) from \eqref{eq:alpha-k}.  They converge to
the continuous, full-support finite measure \(\alpha\).

\begin{lemma}[smooth equalizers with a one-sided logarithmic slope]
\label{lem:equalizers}
There exist positive functions \(q_k\in C^\infty([0,1])\) such that:
\begin{enumerate}
\item \(q_k\equiv1\) in neighborhoods of \(0\) and \(1\);
\item with \(b_k=(k+\kappa_1+1)/4\),
\begin{equation}
 (\log q_k)'(z)\ge-b_k\qquad(0\le z\le1);
 \label{eq:q-slope}
\end{equation}
\item
\begin{equation}
 \sup_{0\le z\le1}
 \abs{\int_0^zq_k(u)\alpha_k(\dd u)-z}
 \longrightarrow0;
 \label{eq:q-equalize}
\end{equation}
\item the construction can be arranged so that
\begin{equation}
 \sup_{z\in[0,1]}|\log q_k(z)|=o(k).
 \label{eq:q-sublinear}
\end{equation}
\end{enumerate}
\end{lemma}

\begin{proof}
Fix a partition depth \(r\), write \(\delta=2^{-r}\), and let
\(I_{r,j}=[j\delta,(j+1)\delta]\), \(0\le j<2^r\).  Since \(\alpha\) has
full support,
\(\alpha(I_{r,j})>0\) for every \(j\).  For fixed \(r\), weak convergence
and atomlessness give
\[
 \alpha_k(I_{r,j})\longrightarrow\alpha(I_{r,j})>0
\]
uniformly over the finite collection of cells.  Define the step equalizer
\begin{equation}
 q_{k,r}^{\mathrm{step}}(z)
 =\frac{\delta}{\alpha_k(I_{r,j})}
 \quad(z\in I_{r,j}^{\circ}).
 \label{eq:step-equalizer}
\end{equation}
It gives exactly mass \(\delta\) to each cell.  Therefore its cumulative
function differs from \(z\) by at most \(\delta\).

For fixed \(r\), all the finitely many values in
\eqref{eq:step-equalizer}, and their reciprocals, remain bounded as
\(k\to\infty\).  Replace the jumps by flat \(C^\infty\) bridges in small
neighborhoods of the dyadic boundaries, interpolate in the logarithmic
variable, and set the function equal to one near the two endpoints.  Upward
bridges have no constraint.  A downward bridge with logarithmic drop \(D\)
can be spread over width slightly larger than \(D/b_k\), ensuring
\eqref{eq:q-slope}.  Since \(b_k\to\infty\), all bridge widths tend to zero
for fixed \(r\).  The measures \(\alpha_k\) converge to an atomless measure,
so the total \(\alpha_k\)-mass of these finitely many shrinking neighborhoods
tends to zero.  The smoothing therefore changes the cumulative integrals by
\(o(1)\), uniformly in the upper endpoint.

Now choose a diagonal sequence.  For each \(r\), take \(K_r\ge r\) so large that
all preceding estimates hold with error at most \(1/r\), the bridges are
disjoint, and the bounded logarithmic size \(C_r\) of the smoothed equalizer
satisfies \(C_r/K_r\le1/r\).  For
\(K_r\le k<K_{r+1}\), use depth \(r\).  Then \(r\to\infty\), so
\(\delta+1/r\to0\), proving \eqref{eq:q-equalize}; the last choice of
\(K_r\) proves \eqref{eq:q-sublinear}.
\end{proof}

The use of dyadic partition boundaries is especially natural here: for
\(r\le k\), every interior partition boundary is a zero of \(P_{k+1}\).
The proof above only needs atomless weak convergence, but these exact zeros make
the smoothing perturbation even smaller.

\subsection{A smooth decreasing exact gauge}

\begin{theorem}[smooth exact Ces\`aro gauge]
\label{thm:smooth-exact}
There exists a positive \(C^\infty\) function \(g_0\), strictly decreasing on
some half-line, such that for every fixed \(t_0\),
\begin{equation}
 \frac1{t-t_0}\int_{t_0}^t\frac{|\Phi(x)|}{g_0(x)}\dd x
 \longrightarrow1.
 \label{eq:smooth-exact-limit}
\end{equation}
It can be chosen with
\begin{equation}
 \log g_0(x)
 =-\frac{(\log x)^2}{2\lambda}-\kappa_1\log x+o(\log x).
 \label{eq:smooth-rough-scale}
\end{equation}
\end{theorem}

\begin{proof}
On the shell \(x=2^k(1+z)\), define
\begin{equation}
 g_0(x)=\frac{g_1(x)}{q_k(z)}.
 \label{eq:g0-shell}
\end{equation}
Because \(q_k\equiv1\) near both shell endpoints, these definitions join with
all derivatives at powers of two; after defining the function arbitrarily on
an initial compact interval, \(g_0\) is globally smooth on a half-line.

Let \(L=k\lambda+\log(1+z)\).  Differentiating \(g_1\) with respect to
\(z\) gives
\begin{equation}
 \frac{\dd}{\dd z}\log g_1(2^k(1+z))
 =-\frac{k+\kappa_1+\log_2(1+z)}{1+z}.
 \label{eq:g1-z-derivative}
\end{equation}
For \(0\le z\le1\), the right side is at most
\(-(k+\kappa_1+1)/2=-2b_k\).  Hence \eqref{eq:q-slope} implies
\[
 \frac{\dd}{\dd z}\log g_0
 =\frac{\dd}{\dd z}\log g_1-(\log q_k)'
 \le -b_k<0
\]
for all sufficiently large \(k\); thus \(g_0\) is strictly decreasing.

By \eqref{eq:normalized-shell} and \eqref{eq:g0-shell},
\[
 \frac{|\Phi(2^k(1+z))|}{g_0(2^k(1+z))}\dd z
 =q_k(z)\alpha_k(\dd z).
\]
Consequently \eqref{eq:q-equalize} gives, uniformly in \(z\),
\begin{equation}
 \int_{2^k}^{2^k(1+z)}\frac{|\Phi(x)|}{g_0(x)}\dd x
 =2^k(z+o(1)).
 \label{eq:equalized-partial-shell}
\end{equation}
In particular a full shell contributes \(2^k(1+o(1))\).  Summing the full
shells before the terminal one and using the geometric Toeplitz argument yields
\[
 \int_{t_0}^{t}\frac{|\Phi(x)|}{g_0(x)}\dd x=t+o(t),
\]
which is \eqref{eq:smooth-exact-limit}.  Finally,
\eqref{eq:q-sublinear} and \(k\asymp\log x\) give
\eqref{eq:smooth-rough-scale}.
\end{proof}

\subsection{Entire approximation preserving monotonicity}

We use the following classical approximation theorem of Hoischen, in the
modern form proved by Gauthier and Kienzle.

\begin{theorem}[Hoischen--Carleman approximation]
\label{thm:Hoischen}
Let \(F\in C^1(\R)\) and let \(\varepsilon:\R\to(0,\infty)\) be continuous.
There is an entire function \(G\), real-valued on \(\R\), such that
\begin{equation}
 |G(x)-F(x)|<\varepsilon(x),
 \qquad
 |G'(x)-F'(x)|<\varepsilon(x)
 \quad(x\in\R).
 \label{eq:Hoischen}
\end{equation}
\end{theorem}

\begin{theorem}[analytic exact Ces\`aro gauge]
\label{thm:analytic-exact}
There exists a zero-free entire function \(g(z)=\exp G(z)\), positive and
strictly decreasing on a real half-line, such that its restriction to that
half-line satisfies \eqref{eq:Q2}.  Moreover it may be chosen so that
\begin{equation}
 \log g(x)
 =-\frac{(\log x)^2}{2\log2}-\kappa_1\log x+o(\log x).
 \label{eq:analytic-rough-scale}
\end{equation}
\end{theorem}

\begin{proof}
Extend \(F=\log g_0\) from \cref{thm:smooth-exact} to a \(C^1\) function on
all of \(\R\).  On a sufficiently large half-line, \(F'(x)<0\).  Choose a
positive continuous error function with
\[
 \varepsilon(x)\longrightarrow0,
 \qquad
 \varepsilon(x)<-\tfrac12F'(x)
\]
on that half-line.  Apply \cref{thm:Hoischen}.  Then
\(G'(x)<F'(x)/2<0\), so \(g(x)=\e^{G(x)}\) is positive and strictly
decreasing.  Also
\[
 \e^{-\varepsilon(x)}
 \le\frac{g(x)}{g_0(x)}\le\e^{\varepsilon(x)}.
\]
Since \(\varepsilon(x)\to0\), splitting the Ces\`aro integral at a fixed
large point and squeezing the tail shows that replacing \(g_0\) by \(g\)
does not change the limit in \eqref{eq:smooth-exact-limit}.  The logarithmic
asymptotic follows from \eqref{eq:smooth-rough-scale} and
\(|G-F|=o(1)\).
\end{proof}

\begin{remark}[constructive status]
The shell equalizers are constructive once a numerical approximation to the
finite measures \(\alpha_k\) is available.  The final entire approximation is
an existence theorem rather than a compact elementary formula.  This is not a
cosmetic defect: \cref{cor:necessary-microstructure} proves that every exact
gauge must acquire unbounded shell distortion.  The simple natural formula
cannot be repaired by a bounded periodic or slowly varying factor.
\end{remark}

\section{A natural exact repair: logarithmic averaging}
\label{sec:logmean}

The additive Ces\`aro mean gives the terminal dyadic shell a fixed positive
fraction of the total interval.  A logarithmic mean gives that shell weight
\(O(1/\log t)\), eliminating the phase profile.

\begin{theorem}[unrestricted logarithmic mean]
\label{thm:logmean}
Let
\begin{equation}
 A_1^{\log}
 =\frac{H_0}{\lambda}\nu\!\left(\frac{W}{1+\cdot}\right)
 =0.09386063575679\ldots\quad\text{(numerical)}.
 \label{eq:Alog}
\end{equation}
Then for every fixed \(t_0>0\),
\begin{equation}
 \boxed{
 \frac1{\log(t/t_0)}\int_{t_0}^t
 \frac{|\Phi(x)|}{A_1^{\log}g_1(x)}\frac{\dd x}{x}
 \longrightarrow1.}
 \label{eq:logmean}
\end{equation}
\end{theorem}

\begin{proof}
The unnormalized mass of the \(k\)-th full logarithmic shell is
\[
 \beta_k
 =\int_{2^k}^{2^{k+1}}\frac{|\Phi(x)|}{g_1(x)}\frac{\dd x}{x}
 =\rho^{-(k+1)}\int_0^1\frac{W(z)}{1+z}P_{k+1}(z)\dd z.
\]
By \eqref{eq:weighted-P-limit},
\(\beta_k\to\beta_\infty=H_0\nu(W/(1+\cdot))\), and the sequence is
bounded.  Up to \(t\in[2^K,2^{K+1})\), the numerator is the sum of
\(K+O(1)\) full shell masses plus a bounded terminal partial mass, while
\(\log(t/t_0)=K\lambda+O(1)\).  Ces\`aro convergence of \((\beta_k)\) gives
\eqref{eq:logmean} with \(A_1^{\log}=\beta_\infty/\lambda\).
\end{proof}

Thus the natural arithmetic formula has two exact interpretations: ordinary
Ces\`aro convergence at dyadic endpoints, and unrestricted convergence for the
multiplicatively invariant logarithmic mean.

\section{Typical rays, RMS scale, and the corrected fluctuation law}
\label{sec:typical}

\subsection{The centered lacunary sum}

Put
\begin{equation}
 \psi(z)=\log|2\sin\pi z|,
 \qquad
 S_n(z)=\sum_{j=0}^{n-1}\psi(\T^jz).
 \label{eq:psi-S}
\end{equation}
Since \(\int_0^1\psi=0\), the Birkhoff sum is centered.  From
\(P_n=2^{-n}\e^{S_n}\) and \eqref{eq:factorization} one obtains, for fixed
\(1<y<2\),
\begin{equation}
 \log|\Phi(2^ky)|
 =-\frac{\log^2(2^ky)}{2\lambda}
  -\kappa_0\log(2^ky)
  +S_{k+1}(y-1)+\Psi(y),
 \label{eq:ray-exact}
\end{equation}
where \(\Psi\) is bounded on compact subintervals of \((1,2)\) and
\begin{equation}
 \kappa_0=\frac32+\log_2\pi.
 \label{eq:kappa0}
\end{equation}
Birkhoff's theorem immediately gives the typical-ray asymptotic.

\begin{theorem}[typical dyadic ray]
\label{thm:typical-ray}
For Lebesgue-almost every \(y\in(1,2)\),
\begin{equation}
 \log|\Phi(2^ky)|
 =-\frac{\log^2(2^ky)}{2\log2}
  -\kappa_0\log(2^ky)+o(k).
 \label{eq:typical-ray}
\end{equation}
\end{theorem}

The normalized remainder does not converge.  Its exact variance and LIL can be
computed without importing a delicate lacunary-series theorem.

\subsection{Exact covariance and variance}

In \(L^2(0,1)\),
\begin{equation}
 \psi(z)=-\sum_{m=1}^{\infty}\frac{\cos(2\pi mz)}{m}.
 \label{eq:psi-fourier}
\end{equation}
Orthogonality gives, for \(r\ge0\),
\begin{equation}
 \int_0^1\psi(z)\psi(2^rz)\dd z
 =\frac{\pi^2}{12\,2^r}.
 \label{eq:covariance}
\end{equation}
Therefore
\begin{equation}
 \boxed{
 \int_0^1S_n(z)^2\dd z
 =\frac{\pi^2}{4}n-\frac{\pi^2}{3}(1-2^{-n}).}
 \label{eq:exact-variance}
\end{equation}
In particular the asymptotic variance per dyadic step is \(\pi^2/4\).

\subsection{Martingale decomposition and LIL}

Let \(P\) be the ordinary Perron operator of the doubling map,
\[
 (P\varphi)(x)=\tfrac12\bigl(\varphi(x/2)+\varphi((x+1)/2)\bigr),
\]
and define
\begin{equation}
 d=2\psi-\psi\circ\T=\log|\tan\pi z|.
 \label{eq:d-martingale}
\end{equation}
The Fourier series shows \(Pd=0\), while
\begin{equation}
 \int_0^1d(z)^2\dd z=\frac{\pi^2}{4}.
 \label{eq:d-variance}
\end{equation}
Moreover the exact telescoping identity is
\begin{equation}
 S_n=\sum_{j=0}^{n-1}d\circ\T^j-\psi+\psi\circ\T^n.
 \label{eq:martingale-decomp}
\end{equation}
The first term is a reverse martingale sum.  The logarithmic singularities of
\(\psi\) have exponential tails, so the two coboundary terms are negligible on
the LIL scale by Borel--Cantelli.

\begin{theorem}[CLT and corrected LIL]
\label{thm:CLT-LIL}
For Lebesgue-almost every \(z\),
\begin{align}
 \frac{S_n}{\sqrt n}&\ \xrightarrow{\ \mathrm d\ }\
 N\!\left(0,\frac{\pi^2}{4}\right),
 \label{eq:CLT}\\
 \limsup_{n\to\infty}\frac{S_n(z)}{\sqrt{2n\log\log n}}
 &=\frac\pi2,
 \qquad
 \liminf_{n\to\infty}\frac{S_n(z)}{\sqrt{2n\log\log n}}
 =-\frac\pi2.
 \label{eq:LIL-standard}
\end{align}
Equivalently, with denominator \(\sqrt{n\log\log n}\), the limsup constant is
\(\pi/\sqrt2\).
\end{theorem}

\begin{proof}
Apply the martingale central limit theorem and the martingale LIL to
\(\sum d\circ\T^j\), using \eqref{eq:d-variance}, then use
\eqref{eq:martingale-decomp}.  The tail estimate
\(\operatorname{Leb}\{|\psi|>u\}=O(\e^{-u})\) makes the coboundary terms
\(o(\sqrt{n\log\log n})\) almost surely.
\end{proof}

Combining \eqref{eq:ray-exact} and \eqref{eq:LIL-standard} yields the physical
frequency form.

\begin{corollary}[physical-ray LIL]
\label{cor:physical-LIL}
For almost every \(y\in(1,2)\),
\begin{equation}
 \limsup_{k\to\infty}
 \frac{\log|\Phi(2^ky)|+\frac{\log^2(2^ky)}{2\lambda}
       +\kappa_0\log(2^ky)}
 {\sqrt{\log(2^ky)\,\log\log\log(2^ky)}}
 =\frac{\pi}{\sqrt{2\lambda}},
 \label{eq:physical-LIL}
\end{equation}
with the negative of this constant as the liminf.
\end{corollary}

\subsection{The \texorpdfstring{\(\sqrt2\)}{sqrt(2)} error in the published lacunary-product constant}

The supplied Aistleitner--Hofer--Larcher source computes its variance by
summing Fourier-coefficient products but omits the factor \(1/2\) in
\(\int_0^1\cos^2(2\pi jx)\dd x=1/2\).  It consequently prints
\(\sigma^2=\pi^2/2\) instead of \(\pi^2/4\).  The exact finite-\(n\) formula
\eqref{eq:exact-variance} and the martingale identity independently force the
smaller value.

\begin{auditfinding}[normalization of the correction]
\label{finding:LIL-correction}
The correct constants are:
\begin{itemize}
\item \(\pi/2\) with denominator \(\sqrt{2n\log\log n}\);
\item \(\pi/\sqrt2\) with denominator \(\sqrt{n\log\log n}\);
\item \(\pi/\sqrt{2\log2}\) in the physical-frequency normalization
\eqref{eq:physical-LIL}.
\end{itemize}
Each is smaller by \(\sqrt2\) than the corresponding printed constant in the
cited paper and in the wave--1 document that copied it.
\end{auditfinding}

\subsection{The exact RMS exponent}

For \(q=2\),
\begin{equation}
 \int_0^1P_n(z)^2\dd z=2^{-n}.
 \label{eq:Pn-L2}
\end{equation}
One proof is the identity \(\Lop_2\mathbf1=\tfrac12\mathbf1\); another is
Fourier orthogonality.  Thus \(\rho_2=1/2\),
\(\Lambda_2=\rho_2^{1/2}=1/\sqrt2\), and the scale dictionary gives
\begin{equation}
 \kappa_2=\kappa(1/\sqrt2)=\log_2(2\pi).
 \label{eq:kappa2}
\end{equation}
This identifies the exponent in the deleted early conjecture: it is the
root-mean-square exponent, not the typical, arithmetic, or extremal exponent.

\section{The sharp extremal product and the coarse uniform envelope}
\label{sec:extremal}

Let \(u_j=\sin^2(\pi2^jz)\), so
\(u_{j+1}=4u_j(1-u_j)\).  The following elementary subaction gives the exact
Gelfond exponent.

\begin{proposition}[sharp Gelfond bound]
\label{prop:gelfond}
For every \(n\ge1\) and \(z\in\R\),
\begin{equation}
 P_n(z)\le\sqrt{\frac53}\left(\frac{\sqrt3}{2}\right)^n.
 \label{eq:gelfond}
\end{equation}
The exponential rate \(\sqrt3/2\) is sharp and is attained at
\(z=1/3,2/3\).
\end{proposition}

\begin{proof}
Set \(V(u)=1+2u/3\).  Direct expansion gives
\begin{equation}
 3V(u)-4uV(4u(1-u))
 =\frac{(2u+1)(4u-3)^2}{3}\ge0.
 \label{eq:subaction-identity}
\end{equation}
Hence
\(u_jV(u_{j+1})\le(3/4)V(u_j)\).  Multiplication yields
\[
 P_n(z)^2V(u_n)\le(3/4)^nV(u_0)
 \le(3/4)^n\frac53,
\]
which is \eqref{eq:gelfond}.  At \(z=1/3\) or \(2/3\), every
\(u_j=3/4\), so \(P_n=(\sqrt3/2)^n\).
\end{proof}

The scale dictionary now gives
\begin{equation}
 \kappa_\infty
 =\kappa(\sqrt3/2)
 =\frac32+\log_2\frac\pi{\sqrt3}.
 \label{eq:kappainfty}
\end{equation}

\begin{theorem}[coarse global envelope]
\label{thm:coarse-envelope}
There is a constant \(C\) such that, for all sufficiently large \(x\),
\begin{equation}
 |\Phi(x)|\le C E(x)x^{-\kappa_\infty}.
 \label{eq:coarse-envelope}
\end{equation}
The power exponent is sharp: on the ray \(x=2^k(4/3)\), the ratio of the two
sides stays between positive constants.
\end{theorem}

\begin{proof}
Insert \eqref{eq:gelfond} into
\eqref{eq:normalized-shell-general} with
\(\Lambda=\sqrt3/2\).  Sharpness follows from the period-two orbit
\(z=1/3\) and positivity of the bounded mantissa factor at \(y=4/3\).
\end{proof}

Consequently
\begin{equation}
 \limsup_{x\to\infty}\frac{\log|\Phi(x)|}{(\log x)^2}
 =-\frac1{2\log2},
 \qquad
 \liminf_{x\to\infty}\frac{\log|\Phi(x)|}{(\log x)^2}=-\infty,
 \label{eq:universal-lims}
\end{equation}
where the second statement follows from the integer zeros.

\section{Largest peak in a dyadic annulus: the Lambert-\texorpdfstring{\(W\)}{W} law}
\label{sec:annular}

The coarse envelope does not yet determine
\[
 M_k=\max_{2^k\le x\le2^{k+1}}|\Phi(x)|.
\]
A fixed maximizing ray such as \(y=4/3\) pays the linear denominator penalty
\((4/3)^{-(k+1)}\).  The true annular maximizer moves toward \(y=1\), paying
only a sublinear boundary-layer loss while spending a slowly increasing number
of doublings reaching the maximizing cycle.

Put
\begin{equation}
 \beta=\log\frac{\sqrt3}{2},
 \qquad
 c=\frac{\lambda}{\pi}\sqrt{\frac32}
 =0.270222319733365\ldots .
 \label{eq:beta-c}
\end{equation}
From \eqref{eq:factorization}, with \(n=k+1\),
\begin{equation}
 M_k=2^{-k(k+1)/2}\pi^{-n}A_n,
 \quad
 A_n=\sup_{0\le\theta\le1}
 H(1+\theta)(1+\theta)^{-n}P_n(\theta).
 \label{eq:M-A}
\end{equation}

\subsection{Boundary-layer reduction}

Write a small \(\theta\) as
\[
 \theta=2^{-r}s,
 \qquad s\in[1/4,1/2].
\]
The first \(r\) factors obey, uniformly in \(s\),
\begin{equation}
 \sum_{j=0}^{r-1}\log|\sin(\pi2^j\theta)|
 =-\frac\lambda2r^2+r\left(\log(\pi s)-\frac\lambda2\right)+O(1).
 \label{eq:small-block}
\end{equation}
The remaining orbit can be chosen to shadow the maximizing period-two cycle.
The exact Gelfond upper bound and a rational-preimage lower construction give
\begin{equation}
 \log P_{n-r}(s)=(n-r)\beta+O((\log\log n)^2)
 \label{eq:tail-max}
\end{equation}
for the optimization, uniformly at the required scale.  The loss comes from a
preimage depth \(O(\log\log n)\).

After removing \(n\beta\), the finite-dimensional objective is
\begin{equation}
 \mathcal F_n(r,s)
 =-\frac\lambda2r^2
 +r\left(\log(\pi s)-\frac\lambda2-\beta\right)
 -n\log(1+s2^{-r}).
 \label{eq:F-objective}
\end{equation}

\subsection{Exact Lambert optimization}

Set
\begin{equation}
 C_0=\log\pi-\frac\lambda2-\beta
 =\log\!\left(\pi\sqrt{\frac23}\right),
 \qquad A=C_0+\log s,
 \qquad z=\lambda r-A.
 \label{eq:C0-z}
\end{equation}
Replacing \(-n\log(1+s2^{-r})\) by \(-ns2^{-r}\) changes the maximum by
\(o(1)\).  A direct calculation gives
\begin{equation}
 \mathcal F_n(r,s)
 =-\frac{z^2}{2\lambda}-n\e^{-C_0}\e^{-z}
 +\frac{A^2}{2\lambda}+o(1).
 \label{eq:F-z}
\end{equation}
As \(s\) ranges through an interval of multiplicative width two, \(A\) ranges
through an interval of length \(\lambda\); successive integer values of \(r\)
therefore tile the real \(z\)-axis.  The bounded \(A^2\)-term is harmless.
The remaining continuous maximum satisfies
\[
 z\e^z=\lambda n\e^{-C_0}=cn,
\]
so \(z=\Wlam(cn)\), the principal Lambert function.

\begin{theorem}[sharp dyadic-annulus maximum]
\label{thm:annular-max}
Let \(w_k=\Wlam(c(k+1))\).  Then
\begin{equation}
 \boxed{
 \begin{aligned}
 \log M_k={}&-\frac\lambda2k(k+1)-(k+1)\log\pi+(k+1)\beta\\
 &-\frac{w_k^2+2w_k}{2\lambda}
 +O((\log\log k)^2).
 \end{aligned}}
 \label{eq:annular-main}
\end{equation}
Equivalently,
\begin{align}
 \log M_k={}&-\frac\lambda2k^2
 +(\beta-\log\pi-\lambda/2)k
 -\frac{(\log k)^2}{2\lambda}
 +\frac{\log k\log\log k}{\lambda}
 \notag\\
 &-\frac{1+\log c}{\lambda}\log k
 +O((\log\log k)^2).
 \label{eq:annular-expanded}
\end{align}
A near-maximizer has
\begin{equation}
 \frac{x}{2^k}=1+\Theta\!\left(\frac{\Wlam(c(k+1))}{k+1}\right),
 \qquad
 r=\frac{\Wlam(c(k+1))}{\lambda}+O(1).
 \label{eq:annular-location}
\end{equation}
\end{theorem}

\begin{proof}
Equations \eqref{eq:small-block}--\eqref{eq:F-objective}, together with the
sharp tail upper bound and the dense-preimage lower construction, reduce
\(\log A_n-n\beta\) to \(\sup\mathcal F_n\) with error
\(O((\log\log n)^2)\).  The tiling argument in
\eqref{eq:C0-z}--\eqref{eq:F-z} gives
\[
 \sup\mathcal F_n
 =-\frac{\Wlam(cn)^2+2\Wlam(cn)}{2\lambda}+O(1).
\]
Insert this into \eqref{eq:M-A}.  Expanding
\(\Wlam(cn)=\log n-\log\log n+\log c+
O(\log\log n/\log n)\) gives \eqref{eq:annular-expanded}.  The critical
point equation also gives
\(n\theta=\Wlam(cn)/\lambda\), proving
\eqref{eq:annular-location}.
\end{proof}

The extra \(- (\log k)^2/(2\lambda)\) term is a second log-normal suppression
in \(\log\log x\).  It cannot be replaced by a fixed power of \(\log x\).

\section{Two-adic exceptional valleys and lobe maxima}
\label{sec:two-adic}

Let \(E_m\) denote the unique maximum of \(|\Phi|\) on \([m,m+1]\).  Large
powers of two create exceptionally deep valleys because many sinc factors
vanish at the left endpoint.

\begin{theorem}[peaks immediately after \(m2^k\)]
\label{thm:two-adic-peaks}
Fix an odd positive integer \(m\), and write the unique maximizer in
\((m2^k,m2^k+1)\) as \(m2^k+t_k\).  Then
\begin{equation}
 (k+1)(1-t_k)\longrightarrow1,
 \label{eq:peak-location}
\end{equation}
and
\begin{equation}
 \boxed{
 E_{m2^k}\sim
 \frac{H(1)|H(m)|}{\e(k+1)}(m2^k)^{-(k+1)}.}
 \label{eq:two-adic-peak}
\end{equation}
In particular,
\begin{equation}
 E_{2^k}\sim\frac{H(1)^2}{\e(k+1)}2^{-k(k+1)},
 \qquad H(1)=0.553771275888810\ldots .
 \label{eq:power-two-peak}
\end{equation}
\end{theorem}

\begin{proof}
Uniformly for \(t\in[0,1]\), the factors that vanish at \(m2^k\) can be
separated to give
\begin{equation}
 |\Phi(m2^k+t)|
 =(m2^k)^{-(k+1)}|H(m)|\,t^{k+1}\Phi(t)(1+o(1)).
 \label{eq:near-two-adic-factor}
\end{equation}
As \(t\uparrow1\),
\(\Phi(t)=H(1)(1-t)(1+o(1))\).  Thus the remaining maximization is the
standard boundary problem
\(t^p(1-t)\), \(p=k+1\), whose maximizer satisfies
\(p(1-t)\to1\) and whose maximum is asymptotic to \(1/(\e p)\).
Strict uniqueness follows from \cref{thm:one-peak}.
\end{proof}

Comparing \eqref{eq:power-two-peak} with \eqref{eq:annular-main} gives
\[
 \log\frac{E_{2^k}}{M_k}=-\frac\lambda2k^2+O(k),
\]
so the first lobe after a power of two is vastly smaller than the largest lobe
in the same dyadic annulus.

\section{Numerical adjudication and reproducibility}
\label{sec:numerics}

All theorem-level conclusions above are independent of the decimal values in
this section.  The numerics are useful for checking normalizations and for
correcting two coarse values in the comparative audit.

\subsection{Spectral and shell constants}

Chebyshev--Lobatto collocation of \(\Lop\), with barycentric evaluation at the
two inverse branches, is stable from 32 to 80 nodes and gives
\begin{align}
 \rho&=0.661322602060565\ldots,\notag\\
 \kappa_1&=2.748070014871334\ldots,\notag\\
 A_1&=0.09126612413159\ldots,\notag\\
 A_1^{\log}&=0.09386063575679\ldots .
 \label{eq:numeric-constants}
\end{align}
An independent direct quadrature on all \(2^{18}\) dyadic cells, using
16-point Gauss--Legendre quadrature in each cell, gives
\[
 A_1^{(18)}=0.09126612413156757,
 \qquad
 (A_1^{\log})^{(18)}=0.09386063575678312.
\]

\subsection{Corrected profile range}

The comparative audit reported a range \([0.9249,1.1133]\), explicitly on a
mantissa grid of spacing \(0.025\).  That grid misses the true narrow
excursions.  The same \(n=18\) cell quadrature, with within-cell root refinement
for \(\mathscr P_n'\), gives
\begin{center}
\begin{tabular}{@{}lll@{}}
\toprule
quantity & value & location\\
\midrule
finite-level minimum
 & \(0.9162685832208\)
 & \(y\approx1.142840221249\)\\
finite-level maximum
 & \(1.1181397445438\)
 & \(y\approx1.428632869969\)\\
\bottomrule
\end{tabular}
\end{center}
At level \(n=19\), the dyadic-cell boundary values are already
\(0.9162686675\) and \(1.1181395810\), supporting convergence to nearby
limits.  These are high-resolution numerical estimates, not certified exact
constants for the limiting singular profile.

\subsection{Minimal reproducibility pseudocode}

\begin{lstlisting}[style=code,caption={Direct dyadic-cell quadrature for the natural profile}]
rho = 0.661322602060565
lam = log(2)
kappa = 0.5 + log(pi/rho)/lam
n = 18
N = 2**n

for each dyadic cell I_j = [j/N,(j+1)/N]:
    use 16-point Gauss-Legendre nodes z in I_j
    logP = sum_{ell=0}^{n-1} log(abs(sin(pi * frac(2**ell*z))))
    logH = sum_{m=1}^{55} log(sin(pi*(1+z)/2**m)/(pi*(1+z)/2**m))
    logW = log(rho/pi) + logH + (kappa-1)*log(1+z)
           + log(1+z)**2/(2*lam)
    density = exp(-n*log(rho) + logW + logP)
    cell_mass[j] = Gauss_integral(density, I_j)

A1 = sum(cell_mass)
F_at_boundaries = cumulative_sum(cell_mass/A1)
profile = (1 + F_at_boundaries)/(1 + arange(N+1)/N)
\end{lstlisting}

The exact variance \eqref{eq:exact-variance} supplies an especially strong
independent check: at \(n=1000\) it gives
\(2464.1112321386\ldots\), asymptotically half the value implied by the
erroneous \(\pi^2/2\) variance.

\section{Document-by-document adjudication}
\label{sec:audit}

\begin{longtable}{@{}p{0.17\textwidth}p{0.36\textwidth}p{0.39\textwidth}@{}}
\toprule
source & strongest correct contribution & correction or limitation\\
\midrule
\endfirsthead
\toprule
source & strongest correct contribution & correction or limitation\\
\midrule
\endhead
Wave--1 document 1
 & Exact dyadic factorization; typical-ray exponent; covariance, variance
 \(\pi^2/4\), and the correct LIL normalization; coarse extremal analysis.
 & Does not solve the arithmetic Ces\`aro problem.  Its annular expansion is
 correct only to a coarser remainder than document 2.\\[0.5em]
Wave--1 document 2
 & Sharp Gelfond subaction; unique lobe peak; full Lambert-\(W\) annular
 asymptotic; two-adic peak theorem.
 & Deliberately follows the pointwise/extremal track and does not determine the
 arithmetic transfer eigenvalue or the original Ces\`aro gauge.\\[0.5em]
Wave--1 document 3
 & Broad synthesis; correct separation of typical, arithmetic, RMS, and
 extremal exponents; construction of the natural analytic gauge satisfying
 the bounded requirement.
 & Inherits the published Aistleitner--Hofer--Larcher LIL constant, which is too
 large by \(\sqrt2\).  Its discussion of failure of the all-time limit applies
 to the natural gauge, not to all possible analytic gauges.\\[0.5em]
Wave--1 document 4
 & Most systematic transfer-operator and \(L^q\) treatment; correct
 arithmetic exponent, RMS identity, singular-profile mechanism, and bounded
 answer.
 & The spectral-gap package is imported rather than fully proved; fluctuation
 theory and the sharp Lambert correction are not developed.  Conditional
 intermediate-rate heuristics should not be read as established theorems.\\[0.5em]
Aistleitner--Hofer--Larcher
 & Identifies the Fouvry--Mauduit arithmetic constant and proves the rigorous
 enclosure \(0.66130<\rho<0.66135\); supplies deep lacunary-product estimates.
 & Its variance computation omits the Parseval factor \(1/2\), printing
 \(\pi^2/2\) instead of \(\pi^2/4\); affected LIL constants must be divided by
 \(\sqrt2\).\\[0.5em]
Wave--2 comparative audit
 & Correctly reconciles the exponent spectrum, proves singularity of the
 arithmetic eigenmeasure, identifies the natural phase profile, and introduces
 the exact logarithmic-mean repair.
 & The reported profile range is a coarse-grid range.  The implication
 ``Hardy-field, hence every elementary formula'' is too broad.  Most
 importantly, ``unattainable in principle'' is false without a bounded- or
 stabilizing-distortion hypothesis: \cref{thm:analytic-exact} gives an
 unrestricted analytic exact gauge.\\
\bottomrule
\end{longtable}

\section{Final theorem and practical interpretation}
\label{sec:final}

The results can be compressed into one theorem that answers every interpretation
appearing in the supplied materials.

\begin{theorem}[complete Fourier-decay resolution]
\label{thm:complete}
Let
\(\Phi(x)=\prod_{n\ge0}\sinc(\pi x/2^n)\).  Then:
\begin{enumerate}
\item There is no positive global pointwise asymptotic equivalent.  The
universal log-square limsup is \(-1/(2\log2)\), while integer zeros force
liminf \(-\infty\).
\item For almost every dyadic mantissa, the power correction is
\(\kappa_0=3/2+\log_2\pi\), with CLT variance \(\pi^2/4\) and physical LIL
constant \(\pi/\sqrt{2\log2}\).
\item The natural arithmetic gauge is
\(g_{\mathrm{nat}}=A_1E(x)x^{-\kappa_1}\), with
\(\kappa_1=1/2+\log_2(\pi/\rho)\).  It satisfies the bounded condition
\eqref{eq:Q3}; its unrestricted limit points form the range of
\(\mathscr P(y)=(1+F(y-1))/y\), and it gives limit one at dyadic endpoints.
\item No gauge with bounded shell distortion relative to \(g_1\), no
stabilizing-shell gauge, and no Hardy-field gauge of the correct scale can
satisfy the exact unrestricted limit \eqref{eq:Q2}.
\item Nevertheless, there exists a positive real-analytic, eventually strictly
decreasing gauge satisfying \eqref{eq:Q2}.  It can have the same rough
log-normal--power scale as \(g_1\), but must develop unbounded fine shell
microstructure.
\item The natural gauge also gives an exact unrestricted logarithmic mean after
multiplication by \(A_1^{\log}\).
\item The RMS power is \(\kappa_2=\log_2(2\pi)\); the coarse uniform power is
\(\kappa_\infty=3/2+\log_2(\pi/\sqrt3)\).
\item The largest dyadic-annulus peak obeys the Lambert-\(W\) formula
\eqref{eq:annular-main}; lobe peaks after \(m2^k\) obey
\eqref{eq:two-adic-peak}.
\end{enumerate}
\end{theorem}

For applications, the choice of answer is therefore straightforward.

\begin{itemize}
\item Use \(A_1g_1\) for a simple analytic arithmetic envelope, two-sided
Ces\`aro control, or dyadic-endpoint normalization.
\item Use \(A_1^{\log}g_1\) when multiplicative/logarithmic averaging is
natural and an unrestricted exact limit is desired.
\item Use the adaptive analytic gauge of \cref{thm:analytic-exact} only when the
literal additive Ces\`aro limit is essential; it is necessarily noncanonical.
\item Use \(\kappa_0\), \(\kappa_2\), or \(\kappa_\infty\) only for the
specific typical, RMS, or extremal notion they encode.
\item Use the Lambert formula, not a fixed logarithmic power, for the largest
peak in a dyadic annulus.
\end{itemize}

\begin{boxedremark}
The phrase ``how fast does the Fourier transform decay?'' has one universal
first answer and several incompatible second answers.  Universally, the
quadratic logarithmic coefficient is \(-1/(2\log2)\).  Beyond that order, the
correct exponent and correction depend on whether one asks for typical values,
means, RMS size, a coarse envelope, an annular maximum, or a specified
exceptional lobe.  The original additive Ces\`aro existence question is yes,
but only through a deliberately adaptive analytic gauge; the simple canonical
gauge has a singular nonconstant phase profile instead.
\end{boxedremark}

\appendix

\section{Details behind the exponent spectrum}
\label{app:spectrum}

For \(q>0\), let \(\rho_q\) be the leading eigenvalue of \(\Lop_q\) and set
\begin{equation}
 \Lambda_q=\rho_q^{1/q},
 \qquad
 \kappa_q=\frac{\log\pi+\lambda/2-\log\Lambda_q}{\lambda}.
 \label{eq:q-spectrum}
\end{equation}
Then the \(q\)-th moment of the normalized shell ratio
\(|\Phi|/(E x^{-\kappa_q})\) has a finite nonzero limit under the corresponding
RPF hypotheses.  At the endpoints:
\begin{itemize}
\item \(q\downarrow0\):
\(\log\Lambda_0=\int_0^1\log|\sin\pi z|\dd z=-\log2\), giving
\(\kappa_0\).
\item \(q=1\): \(\Lambda_1=\rho\), giving \(\kappa_1\).
\item \(q=2\): \(\rho_2=1/2\), hence \(\Lambda_2=1/\sqrt2\), giving
\(\kappa_2\).
\item \(q\to\infty\): \(\Lambda_\infty=\sqrt3/2\), giving
\(\kappa_\infty\).
\end{itemize}
The map \(q\mapsto\log\rho_q\) is convex by H\"older's inequality, which
organizes the monotone interpolation between these distinguished values.

\section{Why the adaptive construction does not contradict singularity}
\label{app:no-contradiction}

For the natural gauge, the normalized shell measures converge to the singular
measure \(\mu\).  A bounded multiplier cannot change the measure class, which
is exactly \cref{thm:bounded-distortion}.  The equalizers
\(q_k\) in \cref{lem:equalizers} do something qualitatively different.  At
partition depth \(r\), they assign reciprocal mass
\[
 q_{k,r}|_{I_{r,j}}\approx\frac{2^{-r}}{\alpha_k(I_{r,j})}.
\]
As \(r\to\infty\), singularity forces these reciprocals to have unbounded
relative variation.  The depth grows only after \(k\) is large enough that the
required downward logarithmic bridges fit inside the cells while respecting
monotonicity of \(g_1/q_k\).  Thus the construction converges weakly to
Lebesgue at every fixed resolution, but has no bounded limiting density with
respect to \(\mu\).  That is precisely the escape route left open by the
singular-profile theorem.

\section{Source map and bibliographic notes}
\label{app:sources}

The repository paths below are included so that every audited claim can be
traced to its source.  The present article does not assume the correctness of
those drafts; the proofs above adjudicate their overlapping statements.

\begin{thebibliography}{99}

\bibitem{AHL}
C.~Aistleitner, R.~Hofer, and G.~Larcher,
\newblock On evil Kronecker sequences and lacunary trigonometric products,
\newblock \emph{Annales de l'Institut Fourier} \textbf{67} (2017), no.~2,
637--687,
\newblock DOI: \href{https://doi.org/10.5802/aif.3094}{10.5802/aif.3094};
preprint arXiv:1502.06738.

\bibitem{FouvryMauduit}
E.~Fouvry and C.~Mauduit,
\newblock Sommes des chiffres et nombres presque premiers,
\newblock \emph{Mathematische Annalen} \textbf{305} (1996), 571--599.

\bibitem{GauthierKienzle}
P.~M. Gauthier and J.~Kienzle,
\newblock Approximation of a function and its derivatives by entire functions,
\newblock \emph{Canadian Mathematical Bulletin} \textbf{59} (2016), 87--94,
\newblock DOI: \href{https://doi.org/10.4153/CMB-2015-060-5}{10.4153/CMB-2015-060-5}.

\bibitem{Hoischen}
L.~Hoischen,
\newblock Approximation und Interpolation durch ganze Funktionen,
\newblock \emph{Journal of Approximation Theory} \textbf{15} (1975), 116--123.

\bibitem{WikipediaFabius}
Wikipedia contributors,
\newblock Fabius function,
\newblock \url{https://en.wikipedia.org/wiki/Fabius_function}, accessed
26 August 2026.

\bibitem{OriginalQuestion}
V.~Reshetnikov,
\newblock Asymptotic decay rate of an infinite product of sinc functions,
\newblock repository transcription of Mathematics Stack Exchange question
2745497,
\newblock \url{https://github.com/VladimirReshetnikov/ProveIt/blob/main/Analysis/FabiusFunction/docs/non-formalized-research-frontiers/drafts/rvachev_up_fourier_decay/asymptotic-decay-rate-of-an-infinite-product-of-sinc-functions/asymptotic-decay-rate-of-an-infinite-product-of-sinc-functions.tex}.

\bibitem{Wave1Doc1}
\newblock Wave--1 document 1,
\newblock \emph{Decay of the Fourier Transform of Rvachev's Up-Function},
\newblock \url{https://github.com/VladimirReshetnikov/ProveIt/blob/main/Analysis/FabiusFunction/docs/non-formalized-research-frontiers/drafts/rvachev_up_fourier_decay/rvachev_up_fourier_decay-1/rvachev_up_fourier_decay.tex}.

\bibitem{Wave1Doc2}
\newblock Wave--1 document 2,
\newblock \emph{Sharp Real-Axis Decay of the Fourier Transform of Rvachev's Up-Function},
\newblock \url{https://github.com/VladimirReshetnikov/ProveIt/blob/main/Analysis/FabiusFunction/docs/non-formalized-research-frontiers/drafts/rvachev_up_fourier_decay/Rvachev_Up_Fourier_Decay-2/Rvachev_Up_Fourier_Decay.tex}.

\bibitem{Wave1Doc3}
\newblock Wave--1 document 3,
\newblock \emph{Sharp Decay Laws for the Fourier Transform of Rvachev's Up-Function},
\newblock \url{https://github.com/VladimirReshetnikov/ProveIt/blob/main/Analysis/FabiusFunction/docs/non-formalized-research-frontiers/drafts/rvachev_up_fourier_decay/Rvachev_Up_Fourier_Decay-3/Fourier_decay_of_Rvachev_up.tex}.

\bibitem{Wave1Doc4}
\newblock Wave--1 document 4,
\newblock \emph{Decay of the Fourier Transform of Rvachev's Up-Function},
\newblock \url{https://github.com/VladimirReshetnikov/ProveIt/blob/main/Analysis/FabiusFunction/docs/non-formalized-research-frontiers/drafts/rvachev_up_fourier_decay/rvachev_up_fourier_decay-4/rvachev_up_fourier_decay.tex}.

\bibitem{ComparativeAudit}
\newblock Wave--2 comparative report,
\newblock \emph{The Decay Envelope of the Rvachev Sinc Product: A Comparative Audit},
\newblock \url{https://github.com/VladimirReshetnikov/ProveIt/blob/main/Analysis/FabiusFunction/docs/non-formalized-research-frontiers/drafts/rvachev_up_fourier_decay/Rvachev_Up_Fourier_Decay_Comparative_Audit/Rvachev_Up_Fourier_Decay_Comparative_Audit.tex}.

\bibitem{RepositoryAHL}
\newblock Repository copy of the Aistleitner--Hofer--Larcher source,
\newblock \url{https://github.com/VladimirReshetnikov/ProveIt/blob/main/Analysis/FabiusFunction/docs/papers/arXiv-1502.06738v2/Lacunary_products.tex}.

\end{thebibliography}

\end{document}
